\section{Discussion}\label{sec:discussion}
\ifnum \workshopsubmission=0
We propose TDS, a practical and asymptotically exact conditional sampling algorithm for diffusion models. 
We compare TDS to other approaches and demonstrate the effectiveness and flexibility of TDS on image class conditional generation and inpainting tasks. 
On protein motif-scaffolding problems with short scaffolds, TDS outperforms the current (conditionally trained) state of the art. 
\fi

A limitation of TDS is its requirement for  additional computes to simulate multiple particles.
While we observe improved performance with just two particles in some cases, the optimal number is problem dependent.
Moreover, the computational efficiency depends on how closely the twisting functions approximate exact conditionals,
which depends on the unconditional model and conditioning information.
Lastly, choosing twisting functions for generic constraints may be challenging.
Our future work will focus on addressing these limitations and improving the computational efficiency. 

%Future works include extensions 
%to large-scale image models and improvements on protein applications. In particular, the latter involves experimenting with a stronger unconditional protein model, and exploring the use of conditionally trained diffusion models as effective proposals for the SMC procedure. 
%Furthermore, there is a potential for exploring TDS's applications to compositional conditional tasks. 